\newpage
\begin{summarybox}
    \textbf{\textcolor{CSummary}{一句话核心}}
    \textbf{同一起点的两条线段若垂直且等长，就把其中一条绕公共端点转 $90^\circ$；坐标先减中心，再交换两个变化量并让其中一个变号，最后加回中心。}

    \medskip
    \textbf{\textcolor{CSummary}{使用条件}}
    \begin{itemize}[leftmargin=*,itemsep=0.5em,topsep=0.4em]
      \item \textbf{条件1（公共端点）：} $B(a,b),\ P(x_1,y_1),\ Q(x_2,y_2)$，比较 $\overrightarrow{BP}$ 与 $\overrightarrow{BQ}$。
      \item \textbf{条件2（两个同时满足）：} $BP\perp BQ$ \textbf{并且} $BP=BQ>0$。
    \end{itemize}

    \medskip
    \textbf{\textcolor{CSummary}{最常用公式}}
    令 $u=x_1-a,\ v=y_1-b$，则
    \[
      \boxed{\begin{aligned}
        \text{逆时针 }90^\circ:&\quad Q=(a-v,b+u),\\
        \text{顺时针 }90^\circ:&\quad Q=(a+v,b-u),\\
        \text{不分方向}:&\quad |x_2-a|=|y_1-b|.
      \end{aligned}}
    \]

    \medskip
    \textbf{\textcolor{CSummary}{考场识别与使用顺序}}
    \begin{itemize}[leftmargin=*,itemsep=0.5em,topsep=0.4em]
      \item \textbf{看到：} 等腰直角三角形、正方形相邻边，或 $BP\perp BQ$ 且 $BP=BQ$。
      \item \textbf{先做：} 写出 $(x_1-a,y_1-b)$，直接求 $Q$ 的两种可能坐标。
      \item \textbf{若 $Q$ 在固定直线：} 先用水平或竖直变化量的绝对值关系，尽量只求一个坐标。
      \item \textbf{后做：} 再代入曲线方程、面积公式或其它条件；避免无目的展开。
    \end{itemize}

    \medskip
    \textbf{\textcolor{CSummary}{关键提醒}}
    \begin{itemize}[leftmargin=*,itemsep=0.5em,topsep=0.4em]
      \item \textbf{不能漏：} 缺少等长不能套旋转等长公式；未知转向必须考虑顺时针与逆时针。
      \item \textbf{不能错：} 不是直接交换 $(x_1,y_1)$，而是交换相对中心的 $(x_1-a,y_1-b)$；两坐标中恰好一个变号。
      \item \textbf{高考母题：} $B(5,0),Q(6,y_2)$ 时立即有 $|y_1|=1$；2020全国III卷椭圆题面积为 $5/2$。
    \end{itemize}
\end{summarybox}
